\rhead{CX206: Programming Methodology Lab}
\phantomsection 
\section*{Programming Methodology Lab (CX206)}
\label{course:CX206}

\noindent \small{\hyperref[main:scheme]{$\leftarrow$ Back to Scheme}} \vspace{0.5cm}

\noindent \textbf{Credits:} 2 (0-0-2)

\subsection*{Course Content}
\noindent\textbf{Lab 1: Imperative vs. Object-Oriented Implementations} \\
Focuses on implementing the same algorithm (e.g., sorting network or matrix operations) using C-style procedural code vs. object-oriented class hierarchies in C++, observing compile-time vs. runtime differences and encapsulation benefits. \\

\vspace{0.2cm}

\noindent\textbf{Lab 2: Cross-Language Function Equivalence} \\
Focuses on writing identical functionality (e.g., data processing pipeline) in C++, Java, and Python, comparing compilation/execution models, type safety, and performance characteristics across the three language ecosystems. \\

\vspace{0.2cm}

\noindent\textbf{Lab 3: Functional Programming Patterns} \\
Focuses on implementing recursive algorithms (e.g., tree traversal, list processing) using functional style with higher-order functions and immutability in Python/JavaScript, contrasting with imperative loops and mutation. \\

\vspace{0.2cm}

\noindent\textbf{Lab 4: Design Patterns Implementation} \\
Focuses on implementing classic design patterns (Strategy, Observer, Factory) in two different languages (C++/Python), comparing syntax differences and runtime behavior for polymorphism and loose coupling. \\

\vspace{0.2cm}

\noindent\textbf{Lab 5: Concurrent Programming Models} \\
Focuses on implementing parallel workloads using threads (C++/Java) vs. goroutines (Go) or async/await (JavaScript/Python), measuring scalability and reasoning about race conditions across paradigms. \\

\vspace{0.2cm}

\noindent\textbf{Lab 6: Domain-Specific Language Usage} \\
Focuses on processing the same dataset using SQL queries, regex patterns, and shell scripting, then integrating results via Python, understanding how DSLs compose with general-purpose languages. \\

\vspace{0.2cm}

\noindent\textbf{Lab 7: Metaprogramming Techniques} \\
Focuses on code generation using C++ templates, Python decorators, and Java annotations, implementing a mini DSL or configuration system that demonstrates compile-time vs. runtime metaprogramming trade-offs. \\

\vspace{0.2cm}

\noindent\textbf{Lab 8: Language Interoperability} \\
Focuses on calling C/C++ code from Python (ctypes/cffi), Java Native Interface, or embedding Python in C++, building a hybrid application that leverages language strengths for different subsystems. \\

\vspace{0.2cm}

\noindent\textbf{Lab 9: Parsing and Lexical Analysis} \\
Focuses on implementing simple lexers/parsers for toy languages using regex, hand-coded state machines, and parser combinator libraries across languages, connecting theory from CS203 to practical language processing. \\

\vspace{0.2cm}

\noindent\textbf{Lab 10: Language Ecosystem Capstone} \\
Focuses on building a polyglot microservice (e.g., data pipeline with Python ETL + Go API + SQL backend), containerizing components, and measuring end-to-end performance across the language/runtime boundaries.

\vspace{0.2cm}

\newpage
\rhead{Curriculum Schema}
